\documentclass[11pt]{amsart}

\usepackage[utf8]{inputenc}
\usepackage{amsmath,amssymb,amsthm,mathtools}
\usepackage[margin=1.05in]{geometry}
\usepackage{enumitem}
\usepackage{microtype}
\usepackage{booktabs}
\usepackage{xcolor}
\usepackage{listings}
\usepackage[colorlinks=true,linkcolor=blue!50!black,citecolor=blue!50!black,urlcolor=blue!50!black]{hyperref}

\lstset{basicstyle=\ttfamily\scriptsize,breaklines=true,columns=fullflexible,
  frame=single,framerule=0.3pt,rulecolor=\color{black!40},showstringspaces=false,
  keepspaces=true,upquote=true}

\newtheorem{theorem}{Theorem}
\renewcommand{\thetheorem}{\Alph{theorem}}
\newtheorem{maincor}[theorem]{Corollary}
\newtheorem{lemma}{Lemma}[section]
\newtheorem{proposition}[lemma]{Proposition}
\newtheorem{corollary}[lemma]{Corollary}
\theoremstyle{definition}
\newtheorem{definition}[lemma]{Definition}
\newtheorem{problem}[lemma]{Problem}
\newtheorem{principle}[lemma]{Principle}
\newtheorem{question}[lemma]{Question}
\theoremstyle{remark}
\newtheorem{remark}[lemma]{Remark}
\newtheorem{example}[lemma]{Example}

\newcommand{\N}{\mathbb{N}}
\newcommand{\Fin}{\mathcal{F}}
\newcommand{\topel}{\operatorname{top}}
\newcommand{\cl}{\prec}

\title[The set-theoretic Yellowstone sequence is a permutation]{The set-theoretic Yellowstone
sequence is a permutation:\\ a top-element descent argument}

\author{Keio}
\thanks{Research draft, October 29, 2026. The results and the exposition were developed
interactively with an AI assistant (Claude, Anthropic); the author is responsible for the
content. The finite data quoted in the paper are reproduced by the script in
Appendix~\ref{app:script}. This draft has not been peer reviewed.}

\date{October 29, 2026}

\begin{document}

\begin{abstract}
The \emph{set-theoretic Yellowstone sequence} (OEIS A252867, F.~T.~Adams-Watters, 2014) is
defined by $a(0)=0$, $a(1)=1$, $a(2)=2$ and, for $n\ge3$, by letting $a(n)$ be the smallest
non-negative integer not already in the sequence whose binary expansion shares at least one
$1$-bit with $a(n-2)$ and no $1$-bit with $a(n-1)$. It is the binary (set) analogue of the
Yellowstone permutation of Applegate, Havermann, Selcoe, Shevelev, Sloane and Zumkeller, and
since 2015 it has been conjectured to be a permutation of the non-negative integers; the
question was posed again as Open Problem 2026.5 in N.~J.~A.~Sloane's June 2026 DIMACS REU
lecture \emph{Interesting unsolved integer sequences}. We prove the conjecture. Equivalently, the squarefree version A252868 is
a permutation of the squarefree positive integers, which settles the conjecture recorded for
that sequence as well.

The proof works for a general class of \emph{colex-greedy intersection sequences}: for every
admissible seed, the greedy rule ``least unused finite set (in colexicographic order) that meets
the set two steps back and avoids the previous set'' eventually uses every non-empty finite set.
The key step is a \emph{top-element descent}: if some element $b_0$ were absent from all late
terms, then $b_0$ together with the largest element of the term two steps back would form an
admissible candidate that is certainly new, so the largest element of a term could never exceed
that of the term two steps earlier; this bounds all late terms, which is impossible for an
injective sequence. (Using an absent element as a certificate of novelty goes back to Nichols's
work on the binary Enots Wolley sequence; what is new is the monotone form, which lets the
argument conclude at once.) Two short exclusion arguments then give every singleton and every finite set.
The proof is elementary and self-contained.
\end{abstract}

\maketitle

\section{Introduction}\label{sec:intro}

\subsection{Background}
Lexicographically earliest sequences defined by local divisibility rules have a reputation for
being easy to compute and hard to analyse. The prototypes are the EKG sequence (A064413),
shown to be a permutation of the positive integers by Lagarias, Rains and
Sloane~\cite{LRS}, and the Yellowstone permutation A098550, in which $a(n)$ is the smallest
unused number that has a common factor with $a(n-2)$ and is coprime to $a(n-1)$; Applegate,
Havermann, Selcoe, Shevelev, Sloane and Zumkeller proved that it is a permutation~\cite{Yell}.
The Enots Wolley sequence A336957 (the rule with the roles of $a(n-1)$ and $a(n-2)$ exchanged)
was shown to contain every integer with at least two distinct prime factors only in
2026~\cite{Nichols26}; its binary analogue A338833 was treated earlier by Nichols~\cite{Nichols22}.

Replacing ``has a common prime factor'' by ``has a common $1$-bit'' turns these rules into rules
about finite subsets of $\N=\{0,1,2,\dots\}$: a non-negative integer $x=\sum_{i\in X}2^i$ is
identified with its \emph{support} $X\subseteq\N$, two integers share a $1$-bit exactly when
their supports intersect, and the usual order of integers becomes the \emph{colexicographic
order} on finite sets. In December 2014 F.~T.~Adams-Watters introduced the binary analogue of the
Yellowstone permutation:
\begin{equation}\label{eq:def}
\begin{gathered}
a(0)=0,\qquad a(1)=1,\qquad a(2)=2,\\
a(n)=\min\bigl\{m\notin\{a(0),\dots,a(n-1)\}:\ m\wedge a(n-2)\neq0,\ m\wedge a(n-1)=0\bigr\}
\quad(n\ge3),
\end{gathered}
\end{equation}
where $\wedge$ is the bitwise AND (OEIS A252867~\cite{A252867}). It begins
\[
0, 1, 2, 5, 10, 4, 3, 12, 17, 6, 9, 18, 8, 7, 24, 33, 14, 32, 11, 36, 19, 40, 16, 13,
48, 15, 80, 34, 20, 35,\dots
\]
The OEIS entry records the conjecture that it is a permutation of the non-negative integers
(comment as modified by N.~J.~A.~Sloane, January~10, 2015), and Sloane later observed that
the graphs of A109812, A252867, A305369 and A305372 ``all have roughly the same, mysterious,
fractal-like structure'' (OEIS A109812, June~3, 2018)~\cite{A252868}.
The question reappears as Open Problem~2026.5 (``Show that A252867 is a permutation of the
positive integers''; since $a(0)=0$, this means that $(a(n))_{n\ge1}$ is a permutation of the
positive integers) in Sloane's June 2026 REU lecture~\cite{REU26}. The
squarefree version A252868, obtained by replacing bit $i$ by the $(i+1)$-st prime, satisfies
$\mathrm{A252868}(n)=\mathrm{A019565}(\mathrm{A252867}(n-1))$, and its entry records the
corresponding conjecture that it is a permutation of the squarefree numbers (``I believe this
is at present only a conjecture'', Sloane, January~10, 2015)~\cite{A252868}.

The integers missing from an initial segment of A252867 are those with many low $1$-bits:
for instance the smallest value missing from the first $4\cdot10^4$ terms is
$12283=10111111111011_2$ (Appendix~\ref{app:script}).
This is the visible reason why a direct, Yellowstone-style counting argument does not obviously
apply.

\subsection{Results}
\begin{theorem}\label{thm:main}
The sequence A252867 is a permutation of the non-negative integers.
\end{theorem}

\begin{maincor}\label{cor:sqf}
The sequence A252868 is a permutation of the squarefree positive integers.
\end{maincor}

Theorem~\ref{thm:main} is the special case of the following statement about general seeds.
For a finite set $X\subseteq\N$ write $\|X\|=\sum_{i\in X}2^i$; the \emph{colexicographic
order} is $X\cl Y\iff\|X\|<\|Y\|$, equivalently $\max(X\triangle Y)\in Y$.

\begin{definition}[Colex-greedy intersection sequences]\label{def:cgi}
Let $X_1,X_2$ be disjoint non-empty finite subsets of $\N$ (the \emph{seed}). Put $S(1)=X_1$,
$S(2)=X_2$, and for $n\ge3$ let $S(n)$ be the colexicographically least finite set
$Y\notin\{S(1),\dots,S(n-1)\}$ with
\begin{equation}\label{eq:rule}
Y\cap S(n-2)\neq\varnothing\qquad\text{and}\qquad Y\cap S(n-1)=\varnothing .
\end{equation}
We call $(S(n))_{n\ge1}$ the \emph{colex-greedy intersection sequence} (CGI sequence) with seed
$(X_1,X_2)$.
\end{definition}

\begin{theorem}\label{thm:general}
For every seed, the CGI sequence is well defined and injective, and every non-empty finite subset
of $\N$ occurs in it. In integer language: for any two positive integers $x_1,x_2$ with
$x_1\wedge x_2=0$, the sequence $a(1)=x_1$, $a(2)=x_2$,
$a(n)=\min\{m\notin\{a(1),\dots,a(n-1)\}: m\wedge a(n-2)\neq0,\ m\wedge a(n-1)=0\}$ is a
permutation of the positive integers.
\end{theorem}

Theorem~\ref{thm:main} follows because the empty set ($=0$) can never satisfy~\eqref{eq:rule},
so A252867 is $0$ followed by the CGI sequence with seed $(\{0\},\{1\})$.

\subsection{The method}
The proof consists of one descent principle and two exclusion steps.

\begin{principle}[Top-element descent]\label{pr:descent}
Suppose an element $b_0$ is absent from all late terms, and let $t(n)$ denote the largest
element of $S(n)$. For every late $n$ the two-element set $\{b_0,t(n-2)\}$ is an admissible
candidate at step $n$ (it meets $S(n-2)$ and avoids $S(n-1)$), and it has never been used (late
terms do not contain $b_0$, early terms are smaller). It lies below $2^{t(n-2)+1}$, so the
minimality of $S(n)$ gives $t(n)\le t(n-2)$. Thus the top elements are non-increasing along each
residue class of $n$ modulo $2$; this bounds all late terms and contradicts injectivity.
\end{principle}

The absent element $b_0$ plays the role of a ``certificate of novelty'': it guarantees that the
candidate $\{b_0,t(n-2)\}$ has never been used, without any knowledge of which large numbers the
sequence has already consumed. This device is not new in itself: Nichols showed that every bit
recurs in the binary Enots Wolley sequence by comparing the first term that contains a much
larger new bit with a smaller candidate containing the absent bit, which certifies that the
candidate is new~\cite[Lemma~1]{Nichols22}. For the Yellowstone-type rule the certificate can be
attached to the top element of the term two steps back, which is automatically admissible, so
the comparison is available at \emph{every} late step and turns into the monotonicity
$t(n)\le t(n-2)$. Once every element is known to recur infinitely often, every
singleton is forced to appear (Lemma~\ref{lem:singletons}), and a missing set $X$ would force
\emph{all} late terms to meet $X$, which is incompatible with the appearance of the infinitely
many singletons outside $X$ (Theorem~\ref{thm:general}). Section~\ref{sec:remarks} explains which
properties of the colexicographic order are used, records what the method does not give, and
lists related open problems.

\section{Basic properties}\label{sec:basic}

Throughout, $(S(n))_{n\ge1}$ is a CGI sequence with seed $(X_1,X_2)$. For a non-empty finite
set $Y$ let $\topel(Y)=\max Y$; then
\begin{equation}\label{eq:topbound}
2^{\topel(Y)}\le\|Y\|<2^{\topel(Y)+1}.
\end{equation}

\begin{lemma}\label{lem:basic}
For every $n\ge3$ the set of admissible candidates in Definition~\ref{def:cgi} is non-empty, so
$S(n)$ is well defined. Moreover:
\begin{enumerate}[label=\textup{(\roman*)},leftmargin=2em]
\item every $S(n)$ is non-empty;
\item $S(n)\cap S(n-1)=\varnothing$ for all $n\ge2$;
\item the map $n\mapsto S(n)$ is injective;
\item for every $K$ there is $N$ with $\|S(n)\|>K$ for all $n\ge N$.
\end{enumerate}
\end{lemma}

\begin{proof}
We argue by induction on $n\ge3$, assuming that $S(1),\dots,S(n-1)$ are defined, non-empty, and
that consecutive ones are disjoint (true for the seed). Pick $d\in S(n-2)$ and an integer $c$
larger than every element of $S(1)\cup\dots\cup S(n-1)$. Then $Y=\{d,c\}$ meets $S(n-2)$; it
avoids $S(n-1)$ because $d\notin S(n-1)$ (consecutive terms are disjoint) and $c$ exceeds all
elements of $S(n-1)$; and $Y$ is unused because $c\in Y$ lies in no earlier term. So the
candidate set is non-empty, and it has a least element because every finite set has only
finitely many colexicographic predecessors (the colexicographic order is a well-order of type
$\omega$); so $S(n)$ exists. It is non-empty because it meets $S(n-2)$, and it is disjoint from
$S(n-1)$ by~\eqref{eq:rule}. This proves the first claim and (i), (ii). Injectivity (iii) holds by
construction for $n\ge3$, and $S(1)\neq S(2)$ because the seeds are disjoint and non-empty;
(iv) follows from (iii) because only finitely many finite sets $Y$ satisfy $\|Y\|\le K$.
\end{proof}

We say that $Y$ is a \emph{candidate at step $n$} if $Y\cap S(n-2)\neq\varnothing$ and
$Y\cap S(n-1)=\varnothing$. The defining property of the sequence will be used in the following
form.

\begin{lemma}[Minimality]\label{lem:min}
If $Y$ is a candidate at step $n\ge3$ and $Y\notin\{S(1),\dots,S(n-1)\}$, then
$\|S(n)\|\le\|Y\|$. Consequently, if a finite set $Y$ never occurs in the sequence, then $Y$ is a
candidate at only finitely many steps.
\end{lemma}

\begin{proof}
The first assertion is the definition of $S(n)$. If $Y$ never occurs and is a candidate at
infinitely many steps $n$, then $\|S(n)\|\le\|Y\|$ and $S(n)\neq Y$, i.e.\ $\|S(n)\|<\|Y\|$, for
infinitely many $n$, contradicting Lemma~\ref{lem:basic}(iii),(iv).
\end{proof}

\section{Top-element descent}\label{sec:descent}

\begin{lemma}\label{lem:recur}
Every $b\in\N$ belongs to $S(n)$ for infinitely many $n$.
\end{lemma}

\begin{proof}
Suppose that $b_0$ belongs to only finitely many terms, and choose $N\ge3$ with $b_0\notin S(n)$
for all $n\ge N$. Put
\[
K:=\max\Bigl(2^{b_0+1},\ \max_{1\le j<N}\|S(j)\|\Bigr).
\]
By Lemma~\ref{lem:basic}(iv) there is $N'\ge N$ such that $\|S(n)\|>2K$ for all $n\ge N'$. For
$n\ge N'$ the top element $t(n):=\topel(S(n))$ satisfies $2^{t(n)+1}>\|S(n)\|>2K$
by~\eqref{eq:topbound}, i.e.\ $2^{t(n)}>K\ge 2^{b_0+1}$; in particular $t(n)>b_0$.

\emph{Claim: $t(n)\le t(n-2)$ for every $n\ge N'+2$.} Consider $Y_n:=\{b_0,t(n-2)\}$. Since
$t(n-2)>b_0$, we have $\|Y_n\|=2^{b_0}+2^{t(n-2)}<2^{t(n-2)+1}$.
\begin{enumerate}[label=\textup{(\alph*)},leftmargin=2em]
\item $Y_n$ is a candidate at step $n$: it meets $S(n-2)$ in $t(n-2)$; it avoids $S(n-1)$ because
$b_0\notin S(n-1)$ (as $n-1\ge N$) and $t(n-2)\notin S(n-1)$ (as $t(n-2)\in S(n-2)$ and
consecutive terms are disjoint, Lemma~\ref{lem:basic}(ii)).
\item $Y_n$ is unused: if $Y_n=S(j)$ with $j<n$, then $j\ge N$ is impossible because $b_0\in Y_n$,
and $j<N$ is impossible because $\|S(j)\|\le K<2^{t(n-2)}\le\|Y_n\|$.
\end{enumerate}
By Lemma~\ref{lem:min}, $\|S(n)\|\le\|Y_n\|<2^{t(n-2)+1}$, and by~\eqref{eq:topbound}
$t(n)\le t(n-2)$. This proves the claim.

By the claim, $t(n)\le E:=\max(t(N'),t(N'+1))$ for all $n\ge N'$ (induction along each residue
class modulo $2$), i.e.\ $\|S(n)\|<2^{E+1}$ for all $n\ge N'$. This contradicts
Lemma~\ref{lem:basic}(iv).
\end{proof}

\begin{remark}
The only properties of the order used in the proof are: the order type is $\omega$
(Lemma~\ref{lem:basic}(iv)), and the size of a set is controlled by its top element as
in~\eqref{eq:topbound}, so that a two-element set $\{b_0,e\}$ with $b_0<e$ precedes every set
with top element $>e$, while sets with large top element lie beyond any given finite initial
segment. The certificate of novelty is the element $b_0$ itself.
\end{remark}

\section{Singletons and arbitrary finite sets}\label{sec:proof}

\begin{lemma}\label{lem:singletons}
Every singleton $\{b\}$, $b\in\N$, occurs in the sequence.
\end{lemma}

\begin{proof}
If $b\in S(n-2)$ for some $n\ge3$, then $b\notin S(n-1)$ by Lemma~\ref{lem:basic}(ii), so
$\{b\}$ is a candidate at step $n$. By Lemma~\ref{lem:recur} this happens for infinitely many
$n$, so $\{b\}$ occurs by Lemma~\ref{lem:min}.
\end{proof}

\begin{proof}[Proof of Theorem~\ref{thm:general}]
Well-definedness and injectivity are Lemma~\ref{lem:basic}. Let $X$ be a non-empty finite set and
suppose that $X$ never occurs. By Lemma~\ref{lem:min} there is $N_1\ge3$ such that $X$ is not a
candidate at any step $n\ge N_1$, that is,
\begin{equation}\label{eq:noncand}
X\cap S(n-2)\neq\varnothing\ \Longrightarrow\ X\cap S(n-1)\neq\varnothing\qquad(n\ge N_1).
\end{equation}
By Lemma~\ref{lem:recur} (applied to any element of $X$) there is $n_0\ge N_1$ with
$X\cap S(n_0)\neq\varnothing$. Applying~\eqref{eq:noncand} with $n=n_0+2,n_0+3,\dots$ gives
inductively $X\cap S(n)\neq\varnothing$ for every $n\ge n_0$. On the other hand, for every
$d\notin X$ the singleton $\{d\}$ occurs (Lemma~\ref{lem:singletons}), say $\{d\}=S(j_d)$, and
$\{d\}\cap X=\varnothing$ forces $j_d<n_0$. Since $X$ is finite there are infinitely many such
$d$, giving infinitely many distinct indices $j_d<n_0$, a contradiction. Hence $X$ occurs. In
integer language, every positive integer occurs, and injectivity gives the permutation statement.
\end{proof}

\begin{proof}[Proof of Theorem~\ref{thm:main} and Corollary~\ref{cor:sqf}]
For $n\ge1$ let $S(n)$ be the support of $\mathrm{A252867}(n)$. Then $S(1)=\{0\}$,
$S(2)=\{1\}$ are disjoint and non-empty, and~\eqref{eq:def} for $n\ge3$ is exactly the rule of
Definition~\ref{def:cgi}: the value $0$ is used at $n=0$ and could never be chosen again anyway,
since it meets no set. Theorem~\ref{thm:general} shows that every positive integer occurs exactly
once among $\mathrm{A252867}(n)$, $n\ge1$; together with $\mathrm{A252867}(0)=0$ this proves
Theorem~\ref{thm:main}. The map $\mathrm{A019565}$, $\sum_{i\in X}2^i\mapsto\prod_{i\in X}p_{i+1}$
($p_j$ the $j$-th prime), is a bijection from $\N$ onto the squarefree positive integers, so
$\mathrm{A252868}(n)=\mathrm{A019565}(\mathrm{A252867}(n-1))$ is a permutation of the squarefree
positive integers.
\end{proof}

\section{Remarks, scope and open problems}\label{sec:remarks}

\subsection{Other ground sets and orders}
Nothing in the argument uses arithmetic. Let $\Omega$ be a countably infinite set listed as
$\omega_0,\omega_1,\dots$ and order its finite subsets colexicographically with respect to this
listing. Then the proofs above apply verbatim and every non-empty finite subset of $\Omega$
occurs in every CGI sequence over $\Omega$. Applied to the primes listed in increasing order,
this gives the squarefree analogue of Theorem~\ref{thm:general} for every admissible seed, with
the squarefree numbers ordered as in A019565 (colexicographically by their sets of prime
factors). The proof does \emph{not} apply to orders in which the size of a set is not dominated
by its top element, for example to squarefree numbers ordered by size (the sequence A252865):
there the candidate $\{b_0,e\}$ corresponds to $p_{b_0+1}p_{e+1}$, which may exceed primes
larger than $p_{e+1}$, so the claim in the proof of Lemma~\ref{lem:recur} breaks down.

\subsection{What the method does not give}
\begin{enumerate}[leftmargin=2em]
\item \emph{Rates.} The proof is ineffective: it gives no bound on the position of a given
integer. The data in Appendix~\ref{app:script} show that $2^b$ first occurs at an index within $7\%$
of $2^b/3$ for $13\le b\le 15$, while integers with many low $1$-bits occur much later.
\item \emph{Longer delays.} For $k\ge3$, start from $k$ pairwise disjoint non-empty finite sets
$S(1),\dots,S(k)$ and let $S(n)$ be the least unused set meeting $S(n-k)$ and avoiding
$S(n-k+1),\dots,S(n-1)$. Then any $k$ consecutive terms are pairwise disjoint, and the
analogues of Lemmas~\ref{lem:recur} and~\ref{lem:singletons} hold with the same proofs (the
descent now reads $t(n)\le t(n-k)$). The final step, however, only shows that the indices of
terms meeting a missing set $X$ have gaps smaller than $k$. Whether all finite sets occur for
$k\ge3$ remains open.
\item \emph{Derived sequences.} A253581 is defined by $\mathrm{A252867}(n)\wedge
\mathrm{A252867}(n+2)$, and its entry conjectures that every positive integer occurs. This is
not implied by Theorem~\ref{thm:main}; the first occurrences are late for integers with many
$1$-bits. For example $447=110111111_2$ does not occur among the values with $n<39998$
(Appendix~\ref{app:script}); it first occurs at $n=65234$ (cf.\ A253603).
\end{enumerate}

\subsection{Related problems}
The ``cup of coffee'' sequence A280864, in which every run of consecutive multiples of a prime
has length exactly two, is conjectured to be a permutation; its OEIS entry records theorems of
Sloane (2017) stating that every prime and every even number appears and that the conjecture is
equivalent to the appearance of every square~\cite{A280864}. The top-element descent does not
apply directly there, because the role of the
forbidden element is played by a prime and the multiplicative order is not top-element
dominated. It would be interesting to find an analogue of Principle~\ref{pr:descent} for
divisibility-defined sequences.

\section*{Acknowledgements}
The author thanks the OEIS contributors whose data and comments made the problem accessible.

\begin{thebibliography}{99}

\bibitem{Yell}
D.~L.~Applegate, H.~Havermann, B.~Selcoe, V.~Shevelev, N.~J.~A.~Sloane and R.~Zumkeller,
\emph{The Yellowstone permutation}, J.\ Integer Seq.\ \textbf{18} (2015), Article 15.6.7;
arXiv:1501.01669.

\bibitem{LRS}
J.~C.~Lagarias, E.~M.~Rains and N.~J.~A.~Sloane, \emph{The EKG sequence}, Experiment.\ Math.\
\textbf{11} (2002), 437--446.

\bibitem{Nichols22}
N.~Nichols, \emph{The binary Enots Wolley sequence}, arXiv:2207.01448 (2022).

\bibitem{Nichols26}
N.~M.~Nichols, \emph{Surjectivity of the Enots Wolley sequence}, arXiv:2609.18054 (2026).

\bibitem{A252867}
OEIS Foundation, \emph{The On-Line Encyclopedia of Integer Sequences}, sequence A252867
(F.~T.~Adams-Watters, Dec.\ 2014), \url{https://oeis.org/A252867}.

\bibitem{A252868}
OEIS Foundation, sequence A252868 (squarefree version of A252867),
\url{https://oeis.org/A252868}; see also A019565, A109812, A252865, A253581, A280864,
A305369, A305372.

\bibitem{A280864}
OEIS Foundation, sequence A280864 (the ``cup of coffee'' sequence), with Sloane's Theorems~1
and~5 in the comments, \url{https://oeis.org/A280864}.

\bibitem{REU26}
N.~J.~A.~Sloane, \emph{Interesting unsolved integer sequences}, slides for the REU seminar talk
``Some interesting unsolved problems concerning number sequences'', DIMACS, Rutgers University,
June 11, 2026, Open Problem 2026.5, \url{https://neilsloane.com/doc/REU2026.pdf}.

\end{thebibliography}

\appendix
\section{Verification script}\label{app:script}
The theorems above are proved for all $n$; the following standard-library Python script only
re-computes the finite data quoted in the paper (agreement with the OEIS data of A252867 and
A252868, disjointness of consecutive terms, first occurrences of powers of two and of all
$m<2^{12}$, and an illustration of Theorem~\ref{thm:general} for several seeds). Its output is
reproduced after the listing.

\begin{lstlisting}[language=Python]
#!/usr/bin/env python3
"""Sanity checks for the paper on A252867 (standard library only).

The theorem in the paper is proved for all n; this script only re-computes the
finite data quoted in the paper:
  (1) the first 68 terms of A252867 agree with the OEIS data;
  (2) consecutive terms are disjoint and a(n) meets a(n-2) (n >= 3);
  (3) first occurrence indices of 2^b (b <= 15) and of every m < 2^12;
  (4) A252868(n) = A019565(A252867(n-1)) reproduces the OEIS data of A252868;
  (5) the same rule started from other admissible seeds also covers all m < 2^10
      within the computed range (illustrating Theorem C);
  (6) the smallest value missing from the first N terms of A252867 and the smallest
      positive integer missing from the corresponding values of A253581
      (A252867(n) AND A252867(n+2)).
"""
import sys

OEIS_A252867 = [0, 1, 2, 5, 10, 4, 3, 12, 17, 6, 9, 18, 8, 7, 24, 33, 14, 32, 11, 36, 19,
    40, 16, 13, 48, 15, 80, 34, 20, 35, 28, 65, 22, 41, 66, 21, 42, 68, 25, 38, 72, 23, 64,
    26, 69, 50, 73, 52, 67, 44, 81, 46, 129, 30, 97, 130, 29, 98, 132, 27, 100, 131, 56, 70,
    49, 74, 37, 82]
OEIS_A252868 = [1, 2, 3, 10, 21, 5, 6, 35, 22, 15, 14, 33, 7, 30, 77, 26, 105, 13, 42, 65,
    66, 91, 11, 70, 143, 210, 187, 39, 55, 78]

def generate(seeds, N, cap):
    """Greedy rule: a(n) = least unused m with m & a(n-2) != 0 and m & a(n-1) == 0."""
    a = list(seeds)
    used = bytearray(cap)
    for v in a:
        used[v] = 1
    low = 1
    for n in range(len(a), N):
        p2, p1 = a[n - 2], a[n - 1]
        while used[low]:
            low += 1
        m = low
        while True:
            if m >= cap:
                raise RuntimeError("cap too small")
            if not used[m] and (m & p2) and not (m & p1):
                break
            m += 1
        a.append(m)
        used[m] = 1
    return a

def a019565(x):
    primes = []
    c = 2
    while len(primes) < 40:
        if all(c % p for p in primes if p * p <= c):
            primes.append(c)
        c += 1
    r, i = 1, 0
    while x:
        if x & 1:
            r *= primes[i]
        x >>= 1
        i += 1
    return r

def main():
    N = 40000
    a = generate([0, 1, 2], N, 1 << 20)
    ok = a[:len(OEIS_A252867)] == OEIS_A252867
    print("(1) first 68 terms agree with OEIS A252867:", ok)
    disj = all(a[n] & a[n - 1] == 0 for n in range(1, N))
    meet = all(a[n] & a[n - 2] != 0 for n in range(3, N))
    print("(2) consecutive terms disjoint:", disj, "; a(n) meets a(n-2) for n>=3:", meet)
    first = {}
    for i, v in enumerate(a):
        first.setdefault(v, i)
    pw = [first.get(1 << b) for b in range(16)]
    print("(3) index of first occurrence of 2^b, b=0..15:", pw)
    miss = [m for m in range(1 << 12) if m not in first]
    print("    all m < 2^12 occur among the first", N, "terms:", not miss,
          "; latest first occurrence:", max(first[m] for m in range(1 << 12)) if not miss else None)
    b868 = [a019565(a[n]) for n in range(len(OEIS_A252868))]
    print("(4) A019565(A252867) reproduces OEIS A252868:", b868 == OEIS_A252868)
    for seeds in ([1, 2], [3, 4], [6, 1], [5, 10], [1, 6]):
        s = generate(seeds, 30000, 1 << 20)
        fs = set(s)
        allin = all(m in fs for m in range(1, 1 << 10))
        print("(5) seeds", seeds, ": all 1 <= m < 2^10 occur in 30000 terms:", allin)
    smallest_missing = min(m for m in range(1 << 20) if m not in first)
    print("(6) smallest value missing from the first", N, "terms of A252867:",
          smallest_missing, "=", bin(smallest_missing))
    andvals = set(a[n] & a[n + 2] for n in range(1, N - 2))
    print("    smallest positive integer missing from A252867(n) AND A252867(n+2), n < N-2:",
          min(m for m in range(1, 1 << 20) if m not in andvals))
    return 0 if (ok and disj and meet and not miss and b868 == OEIS_A252868) else 1

if __name__ == "__main__":
    sys.exit(main())
\end{lstlisting}

\noindent Output:
\begin{lstlisting}
(1) first 68 terms agree with OEIS A252867: True
(2) consecutive terms disjoint: True ; a(n) meets a(n-2) for n>=3: True
(3) index of first occurrence of 2^b, b=0..15: [1, 2, 5, 12, 22, 17, 42, 74, 114, 298, 439, 823, 1915, 2901, 5680, 11014]
    all m < 2^12 occur among the first 40000 terms: True ; latest first occurrence: 19858
(4) A019565(A252867) reproduces OEIS A252868: True
(5) seeds [1, 2] : all 1 <= m < 2^10 occur in 30000 terms: True
(5) seeds [3, 4] : all 1 <= m < 2^10 occur in 30000 terms: True
(5) seeds [6, 1] : all 1 <= m < 2^10 occur in 30000 terms: True
(5) seeds [5, 10] : all 1 <= m < 2^10 occur in 30000 terms: True
(5) seeds [1, 6] : all 1 <= m < 2^10 occur in 30000 terms: True
(6) smallest value missing from the first 40000 terms of A252867: 12283 = 0b10111111111011
    smallest positive integer missing from A252867(n) AND A252867(n+2), n < N-2: 447
\end{lstlisting}

\end{document}
